% Converted from README.md
\paragraph{Delimično preklapanje izvorne i ciljne memorije}\label{delimiux10dno-preklapanje-izvorne-i-ciljne-memorije}

Program \texttt{05\_overlap.c} prepisuje nisku u lokaciju koja se preklapa sa onom odakle se prepisuje. Ako unesemo reč \texttt{Da}, neće se detektovati gredška, jer prilikom izvrdšavanja za unetu reč neće biti preklapanja. Ponovo pokrenimo program i unesimo reč \texttt{dobro}. Memcheck detektuje problem preklapanja memorijskih lokacija sa koje se kopira i one na koju se kopira.

\begin{minted}{text}
==16178== Source and destination overlap in memcpy(0x5204041, 0x5204040, 5)
==16178== at 0x4C32513: memcpy@@GLIBC_2.14 (in /usr/lib/valgrind/vgpreload_memcheck-amd64-linux.so)
==16178== by 0x400617: main (05_overlap.c:9)
\end{minted}

Ukoliko se unese kraća reč od 3 slova ovo će biti jedina gredška koju imamo. Ukoliko unesemo dužu reč na ulazu, \texttt{printf} će nam prijaviti pristup neinicilizovanoj vrednosti.
